\section{Çalışmanın Amacı}

Yıldırım, dünya genelinde her yıl binlerce ölüme ve milyarlarca dolarlık maddi hasara neden olan önemli bir doğal afettir. Özellikle dağlık bölgelerde yıldırım olayları daha sık gözlemlenmekte ve bu bölgelerde yaşayan insanlar için ciddi bir tehdit oluşturmaktadır.

Bu çalışmanın temel amacı, Hakkari bölgesinde günlük meteorolojik veriler kullanılarak yıldırım olaylarının önceden tahmin edilmesidir. Bu amaç doğrultusunda:

\begin{itemize}
    \item Günlük meteorolojik veriler (sıcaklık, nem, basınç, rüzgar, bulutluluk, yağış) toplanmış ve analiz edilmiştir.
    \item Yıldırım tespit sistemi verileri ile meteorolojik veriler eşleştirilmiştir.
    \item Farklı makine öğrenmesi algoritmaları kullanılarak ikili sınıflandırma modelleri geliştirilmiştir.
    \item Modellerin performansları karşılaştırılarak en uygun model belirlenmiştir.
\end{itemize}

\section{Çalışmanın Önemi}

Yıldırım tahmini, birçok sektör için kritik öneme sahiptir:

\begin{enumerate}
    \item \textbf{Havacılık:} Uçuş güvenliği için yıldırım uyarıları gereklidir.
    \item \textbf{Enerji:} Elektrik hatları ve trafo merkezlerinin korunması.
    \item \textbf{Tarım:} Açık alanda çalışan işçilerin güvenliği.
    \item \textbf{Turizm ve Spor:} Dağcılık, golf gibi açık hava aktiviteleri.
\end{enumerate}

Hakkari bölgesi, yüksek rakımı ve dağlık yapısı nedeniyle yıldırım olaylarının sık yaşandığı bir bölgedir. Bu nedenle bölge için güvenilir bir yıldırım tahmin sistemi geliştirmek büyük önem taşımaktadır.

\section{Çalışmanın Kapsamı}

Bu tez çalışması aşağıdaki kapsamda gerçekleştirilmiştir:

\begin{itemize}
    \item \textbf{Coğrafi Kapsam:} Hakkari ili ve çevresi (6 meteoroloji istasyonu)
    \item \textbf{Zaman Kapsamı:} 2015-2025 yılları arası
    \item \textbf{Veri Kaynakları:} Meteoroloji Genel Müdürlüğü ve Yıldırım Tespit Sistemi
\end{itemize}

\section{Tezin Organizasyonu}

Bu tez altı bölümden oluşmaktadır:

\begin{description}
    \item[Bölüm 1:] Giriş ve çalışmanın amacı
    \item[Bölüm 2:] Literatür taraması ve ilgili çalışmalar
    \item[Bölüm 3:] Kullanılan yöntemler ve algoritmalar
    \item[Bölüm 4:] Veri seti açıklaması ve ön işleme adımları
    \item[Bölüm 5:] Deneysel sonuçlar ve model karşılaştırmaları
    \item[Bölüm 6:] Tartışma, sonuç ve gelecek çalışmalar
\end{description}
